% Theory behind the final report: the statistics we use, with worked examples.
% Build from this folder:  pdflatex Theory.tex
\documentclass[11pt,a4paper]{article}
\usepackage[a4paper,margin=2.2cm]{geometry}
\usepackage[T1]{fontenc}
\usepackage[sc]{mathpazo}
\linespread{1.05}
\usepackage{courier}
\usepackage[scaled=0.9]{helvet}
\usepackage{microtype}
\usepackage{amsmath,amssymb}
\usepackage{booktabs}
\usepackage{tabularx}
\usepackage{enumitem}
\usepackage{xcolor}
\usepackage[hidelinks]{hyperref}
\usepackage{titlesec}
\usepackage{tcolorbox}

\titleformat{\section}{\large\sffamily\bfseries}{\thesection.}{0.5em}{}
\titleformat{\subsection}{\normalsize\sffamily\bfseries}{\thesubsection}{0.5em}{}
\setlist{itemsep=2pt,topsep=3pt}
\setlength{\parindent}{0pt}
\setlength{\parskip}{5pt}

% worked examples go in a light box
\newtcolorbox{example}[1]{colback=gray!6, colframe=gray!45, boxrule=0.5pt, arc=2pt,
  left=6pt, right=6pt, top=4pt, bottom=4pt, fonttitle=\sffamily\bfseries\small,
  coltitle=black, colbacktitle=gray!20, title={Worked example: #1}}

\begin{document}

\begin{center}
{\LARGE\bfseries Theory: the statistics behind the report}\\[4pt]
The Perks of Being a First Mover \quad$\cdot$\quad CS215 Open Project \quad$\cdot$\quad Group 02
\end{center}

This document explains every statistical method in our final report, why it fits the question,
and how to compute it, with a small worked example for each. The examples use small made-up
numbers so they can be checked by hand. The real data has 160 million games.

\begin{center}
\small
\begin{tabularx}{\linewidth}{@{}XXc@{}}
\toprule
Question & Method & Section \\
\midrule
What do $h$, $\hat h$, SE, \dots\ mean? & Definitions & \ref{sec:defs} \\
How big is White's advantage? & Maximum likelihood estimate of $h$ & \ref{sec:mle} \\
How sure are we? & Fisher information, 95\% confidence interval & \ref{sec:ci} \\
Do aggressive openings, or slower time controls, have a larger $h$? & $z$-test for a difference of two estimates & \ref{sec:compare} \\
What does ``significant'' mean with 160M games? & Effect sizes over $p$-values & \ref{sec:bign} \\
\bottomrule
\end{tabularx}
\end{center}

\section{Definitions}\label{sec:defs}

\subsection{Symbols}
\begin{center}
\small
\begin{tabularx}{\linewidth}{@{}lX@{}}
\toprule
Symbol & Meaning \\
\midrule
$R^{\text{White}}_i,\ R^{\text{Black}}_i$ & The two players' Lichess ratings in game $i$. \\
$d_i$ & Rating gap in game $i$: $d_i = R^{\text{White}}_i - R^{\text{Black}}_i$ (positive when White is higher rated). \\
$S_i$ & White's score in game $i$: 1 for a win, $\tfrac12$ for a draw, 0 for a loss. \\
$E(d)$ & Expected score of a player rated $d$ points above the opponent (the Elo formula). \\
$h$ & The \textbf{true} first-move advantage: how many rating points moving first is worth, for the
  games we are describing (all of Lichess, one opening, one speed, \dots). One fixed number that
  nobody knows. \\
$\hat h$ & Our \textbf{estimate} of $h$, worked out from the games we actually have
  (``$h$-hat''). The hat always means ``estimated from data''. \\
$\mathrm{SE}(\hat h)$ & Standard error: how much $\hat h$ would typically change if we repeated the
  study with a fresh batch of games (Section~\ref{sec:ci}). \\
95\% CI & Confidence interval $\hat h \pm 1.96\,\mathrm{SE}$: a range that contains the true $h$ in
  95\% of repeated studies. \\
$n_d$ & Number of games with rating gap $d$. \\
$K$ & $\ln 10 / 400 \approx 0.005757$, which turns $10^{x/400}$ into $e^{Kx}$. \\
$I(\hat h)$ & Fisher information: how sharply the log-likelihood peaks at $\hat h$. $\mathrm{SE} = 1/\sqrt{I}$. \\
$z$ & An estimate divided by its standard error. Close to standard normal if the true value is 0. \\
$p$ & $p$-value: the chance of a $|z|$ at least this big if the true value were 0. \\
\bottomrule
\end{tabularx}
\end{center}

\subsection{$h$ and $\hat h$: why a single number has a distribution}
Once we have our games, $\hat h$ is just one number: put the games in, get (say) $28.2$ out. So
what does it mean to talk about the \emph{distribution} of $\hat h$, its standard error, or a
confidence interval?

The point is that the games themselves are partly luck: a blunder here, a lost-on-time there.
Our games are one random sample of all the games that could have been played. A different month,
or a different 100 games, would give a slightly different $\hat h$. So:
\begin{itemize}
\item $h$ is \textbf{fixed} but unknown: a property of how chess is played on Lichess.
\item $\hat h$ is \textbf{random before we collect the data}, because it's computed from random
  games. After we collect them it's one number, but it's one draw from that randomness.
\item The \textbf{distribution of $\hat h$} (its \emph{sampling distribution}) describes the
  values we \emph{would} get if we repeated the whole study many times with fresh games. Its
  centre is the true $h$, and its spread is the standard error.
\end{itemize}

\textbf{A coin example.} Flip a coin 100 times and get 54 heads. Your estimate of the chance of
heads is $\hat p = 0.54$, one number. But you know another 100 flips might give 49 or 57: the
estimate has a spread of about $\sqrt{0.5 \times 0.5 / 100} = 0.05$, even though the coin's true
chance $p$ never changes. $\hat h$ is exactly like $\hat p$: the 0.54 is our one run, and the
spread describes the other runs we could have had.

\textbf{Where the spread comes from.} The likelihood (Section~\ref{sec:mle}) scores every possible
value of $h$ by how well it explains our games. If the peak is \emph{sharp}, only values very
close to $\hat h$ explain the data, so another batch of games would give nearly the same
estimate: a small spread. If the peak is \emph{flat}, many values fit almost equally well, so
another batch could land far away: a large spread. The Fisher information $I$ measures how sharp
the peak is, and $\mathrm{SE} = 1/\sqrt{I}$. A standard theorem says that with many games,
\[
\hat h \;\approx\; \mathcal N\big(h,\ \mathrm{SE}^2\big),
\]
a bell curve centred on the true $h$.

\textbf{Why it matters for us.} With 100 games, $\mathrm{SE} \approx 35$ points
(Section~\ref{sec:ci}), so another 100 games could easily give $-10$ or $+60$. With 160 million
games it's a few hundredths of a point, so any other month would give almost the same $\hat h$.
And it's what lets us compare: saying ``the aggressive openings have a larger $h$'' only means
something if the gap between the two $\hat h$ values is bigger than this luck-of-the-draw spread.
That's what the $\pm$ and the $z$-tests check.

In the rest of this document, and in the report, numbers like ``$h = +11.9 \pm 0.1$'' are
estimates: strictly they are $\hat h$, with their 95\% CI.

\section{The model: White's advantage in rating points}

\subsection{The Elo expected score}
Lichess ratings follow the Elo idea: if player A is rated $d = R_A - R_B$ points above player B,
A's expected score (win = 1, draw = $\frac12$, loss = 0) is
\[
E(d) = \frac{1}{1 + 10^{-d/400}}
\]
$E(0) = 0.5$, $E(50) = 0.5715$, $E(-50) = 0.4285$. Every 400 points multiply the odds of
winning by 10.

\subsection{Why we don't just compare win rates}
A raw White score mixes the first move with the rating gap. Suppose a 1540 plays a 1570 with
White and wins. The ratings alone expect White to score $E(-30) = 0.4569$, so this win is
$1 - 0.4569 = 0.543$ points ``better than expected''. A win between two 1570s is only $0.5$
better than expected. Comparing every result to what the ratings predict makes all games
comparable, whatever the gap.

\subsection{Adding a White bonus $h$}
We model the first move as a bonus of $h$ rating points for White:
\[
E_i(h) = \frac{1}{1 + 10^{-(d_i + h)/400}}, \qquad d_i = R^{\text{White}}_i - R^{\text{Black}}_i
\]
$h$ is the quantity we estimate. Measuring it in rating points means ``White is worth about
$h$ points'' can be compared across ratings, speeds and openings. Draws are counted as half a
point, and we report draw rates separately.

\section{Estimating $h$: maximum likelihood}\label{sec:mle}

\subsection{The likelihood}
Treat each game's score $S_i$ as an outcome with mean $E_i(h)$. The (Bernoulli) log-likelihood is
\[
\ell(h) = \sum_i \Big[ S_i \ln E_i(h) + (1 - S_i) \ln\big(1 - E_i(h)\big) \Big]
\]
With draws ($S_i = \frac12$) this is a \emph{quasi}-likelihood: it isn't the exact probability
of the data, but it is maximised at the same kind of estimate and is the standard way to handle
fractional scores.

\subsection{Setting the derivative to zero}
Write $K = \ln 10 / 400$, so $E_i = 1/(1 + e^{-K(d_i + h)})$. The logistic function has the
handy derivative
\[
\frac{dE_i}{dh} = K\, E_i (1 - E_i)
\]
Differentiating $\ell$ and simplifying:
\[
\ell'(h) = \sum_i \left[ \frac{S_i}{E_i} - \frac{1 - S_i}{1 - E_i} \right] K E_i (1 - E_i)
         = K \sum_i \big( S_i - E_i(h) \big)
\]
So the maximum-likelihood estimate $\hat h$ is the value where
\[
\boxed{\ \sum_i E_i(\hat h) = \sum_i S_i\ }
\]
that is, \emph{White's expected total score equals White's actual total score}. This is
exactly the intuition: choose the bonus that makes the model agree with what happened.

\subsection{Solving it: Newton's method}
There is no closed form, so we iterate. The second derivative is
\[
\ell''(h) = -K^2 \sum_i E_i (1 - E_i) \;<\; 0,
\]
so $\ell$ is concave and has a single maximum. Newton's method moves to where a straight-line
approximation of $\ell'$ crosses zero:
\[
h_{\text{new}} = h - \frac{\ell'(h)}{\ell''(h)}
= h + \frac{K \sum_i (S_i - E_i)}{K^2 \sum_i E_i (1 - E_i)}
\]
Games with the same rating gap $d$ share the same $E$, so we only need, for each gap $d$, the
number of games $n_d$ and White's total score $S_d$. That is why our program stores counts by
rating gap rather than individual games.

\begin{example}{estimating $h$ from 100 games}
Three groups of games, by rating gap $d$:
\begin{center}
\small
\begin{tabular}{@{}rrrr@{}}
\toprule
Gap $d$ & Games $n_d$ & White's score $S_d$ & Expected at $h=0$: $n_d E(d)$ \\
\midrule
$0$ & 40 & 22 & $40 \times 0.5000 = 20.00$ \\
$+50$ & 30 & 19 & $30 \times 0.5715 = 17.15$ \\
$-50$ & 30 & 13 & $30 \times 0.4285 = 12.85$ \\
\midrule
Total & 100 & 54 & 50.00 \\
\bottomrule
\end{tabular}
\end{center}
White scored 54 where the ratings expect 50, so $h > 0$. Here $K = \ln 10/400 = 0.0057565$.

\textbf{Step 1} ($h = 0$):
\begin{align*}
\ell' &= K\,(54 - 50) = 0.023026, \\
-\ell'' &= K^2 \textstyle\sum n_d E_d(1-E_d) = K^2 \big[40(0.25) + 2 \times 30(0.5715)(0.4285)\big] \\
        &= K^2 \times 24.693 = 0.00081827, \\
\text{step} &= 0.023026 / 0.00081827 = 28.14, \quad\text{giving } h = 28.14
\end{align*}

\textbf{Step 2} ($h = 28.14$): the expected scores are now $0.5404$, $0.6106$ and $0.4686$ per
game, adding up to $53.992$. $\ell' = K(54 - 53.992) = 0.000048$, and $-\ell'' = 0.00081312$.
Step $= 0.06$, giving $h = 28.20$.

\textbf{Step 3}: the expected total is now $54.000$, so the step is 0 and we stop.

$\boxed{\hat h = 28.2 \text{ rating points}}$. Newton converges this fast in practice: our program
needs 3 to 5 steps per fit.
\end{example}

\section{How sure are we: standard error and confidence interval}\label{sec:ci}

\subsection{Fisher information}
How sharply the log-likelihood peaks tells us how precisely the data pins down $h$. The
\emph{observed Fisher information} is the curvature at the maximum,
\[
I(\hat h) = -\ell''(\hat h) = K^2 \sum_i E_i (1 - E_i)
\]
A standard result for maximum-likelihood estimates is that, for large samples, $\hat h$ is
approximately normal around the true $h$ with variance $1/I$:\footnote{The term $E(1-E)$ in $I$ is
the variance of a game that is either won or lost. Draws make results less spread out, so the true
variance is a bit smaller and our intervals are slightly conservative (a little too wide). That
errs on the safe side.}
\[
\hat h \;\approx\; \mathcal N\!\left(h,\ \frac{1}{I(\hat h)}\right), \qquad
\mathrm{SE}(\hat h) = \frac{1}{\sqrt{I(\hat h)}}
\]

\subsection{The 95\% confidence interval}
Since 95\% of a normal distribution lies within $1.96$ standard deviations of its mean,
\[
\hat h \pm 1.96\, \mathrm{SE}(\hat h)
\]
is an approximate 95\% confidence interval: if we repeated the study many times, about 95\% of
such intervals would contain the true $h$. We use the normal distribution and not a
$t$-distribution because the variance comes from the model, not from a small sample estimate,
and with millions of games the two would be identical anyway.

\begin{example}{the confidence interval for the 100 games}
At $\hat h = 28.2$: $I = 0.00081310$, so $\mathrm{SE} = 1/\sqrt{0.00081310} = 35.07$ points.
The 95\% CI is $28.2 \pm 1.96 \times 35.07 = 28.2 \pm 68.7$, that is from $-40.5$ to $+96.9$.

With only 100 games we can't even say whether White has an advantage: the interval includes 0.
Each game adds at most $K^2/4$ to $I$, so the SE shrinks like $1/\sqrt N$. With $N$ games,
$\mathrm{SE} \ge 2/(K\sqrt N) \approx 347/\sqrt N$ points: about 35 for 100 games, 0.35 for a
million, and a few hundredths of a point for our 160 million.
\end{example}

\section{Aggressive vs drawish: the $z$-test for a difference}\label{sec:compare}

\subsection{The test}
We fit one $h$ for all games in the four aggressive openings, $\hat h_A$ with standard error
$\mathrm{SE}_A$, and one for the four drawish openings, $\hat h_D$ with $\mathrm{SE}_D$. They come
from different games, so they are independent, and for independent quantities variances add:
\[
\mathrm{Var}(\hat h_A - \hat h_D) = \mathrm{SE}_A^2 + \mathrm{SE}_D^2,
\qquad
\mathrm{SE}(\hat h_A - \hat h_D) = \sqrt{\mathrm{SE}_A^2 + \mathrm{SE}_D^2}
\]
If the openings made no difference, the true difference would be 0, and
\[
z = \frac{\hat h_A - \hat h_D}{\sqrt{\mathrm{SE}_A^2 + \mathrm{SE}_D^2}}
\]
would be approximately standard normal. The 95\% CI for the difference is
\[
(\hat h_A - \hat h_D) \pm 1.96 \sqrt{\mathrm{SE}_A^2 + \mathrm{SE}_D^2}
\]

\begin{example}{two groups}
Suppose the aggressive openings give $\hat h_A = 30$ with $\mathrm{SE}_A = 2$, and the drawish
openings give $\hat h_D = 22$ with $\mathrm{SE}_D = 3$.

Difference: $30 - 22 = 8$ points. Standard error: $\sqrt{2^2 + 3^2} = \sqrt{13} = 3.61$.

$z = 8 / 3.61 = 2.22$, so $p = 2\,(1 - \Phi(2.22)) = 0.027$, and the 95\% CI is
$8 \pm 1.96 \times 3.61 = 8 \pm 7.1$, that is from $0.9$ to $15.1$ points.

Notice that the individual CIs ($30 \pm 3.9$ and $22 \pm 5.9$) overlap, yet the difference is
significant. Overlapping intervals do \emph{not} mean ``no difference'': always test the
difference itself.
\end{example}

\textbf{The same test for time controls.} Each speed (Bullet, Blitz, Rapid, \dots) is also
fitted as its own set of games, and we compare each speed with the next faster one with exactly
this test, for example $z = (\hat h_{\text{Blitz}} - \hat h_{\text{Bullet}}) /
\sqrt{\mathrm{SE}_{\text{Blitz}}^2 + \mathrm{SE}_{\text{Bullet}}^2}$.

\textbf{What this test doesn't do.} Gambits are played more at some ratings than others, and $h$
depends on rating. So part of a difference between the groups can come from \emph{who plays
them} rather than the openings themselves. The report shows each group's average rating, and a
chart comparing the two groups bracket by bracket, so this can be checked by eye.

\section{Big data: significance vs size}\label{sec:bign}
With 160 million games, standard errors are tiny. A difference of half a rating point can
have $p < 10^{-10}$ and still be completely unimportant for anyone playing chess. So in the
report:
\begin{itemize}
\item every effect is given with its \emph{size} and a 95\% CI (in rating points, or
  percentage points for draw rates), not just a $p$-value;
\item we treat a difference of a few rating points as small, whatever its $p$-value;
\item we read how $h$ changes with rating from the bracket-by-bracket charts.
\end{itemize}

\subsection*{Assumptions to keep in mind}
\begin{itemize}
\item \textbf{Independence.} The formulas assume games are independent. Active players
  contribute many games each, so the real uncertainty is somewhat larger than the CIs.
\item \textbf{The Elo curve.} We assume Lichess ratings predict results through the Elo
  formula. Lichess uses Glicko-2, whose expected score is the same curve for established
  players, so this is a good approximation.
\item \textbf{Large-sample normality.} The CI and $z$-tests rely on $\hat h$ being
  approximately normal, which holds when each fit uses many games. That's why brackets with
  fewer than 1{,}000 games are left out.
\end{itemize}

\end{document}
